\subsection{Structures which do not require a comparison theorem}

\begin{result}{C.1}{Proved}{Strict composition already gives a coherent enriched category}
For a small relative category $(M,\Sigma)$, the collection
\[
 \mathcal G(a,b)=N\Ch_\Sigma(M;a,b),\qquad a,b\in\operatorname{Ob}M,
\]
has strictly associative and strictly unital simplicial composition. The unit is the empty path. The same is true with $\Sigma$-valued vertical components. Applying $\operatorname{Ex}^{\infty}$ to each mapping simplicial set preserves these equations and makes each mapping object Kan. Consequently each construction determines a category enriched in spaces, and hence a Segal presentation. This statement does not identify the enrichment with $L(M,\Sigma)$.
\end{result}
\begin{proof}
For $a,b,c\in M$, concatenation of words, interval diagrams and natural transformations defines a functor
\[
 *:\Ch_\Sigma(M;b,c)\times\Ch_\Sigma(M;a,b)
       \longrightarrow\Ch_\Sigma(M;a,c),
\]
where the first factor is composed after the second. The natural transformations agree at the glued vertex because both components there are $1_b$. The concatenated type map is again endpoint preserving and nondecreasing. The only new identification is the common vertex; no new backward arrow is created. Concatenation of finite words is associative as an equality, including the reindexing convention in Part A. The empty word changes neither the diagram nor any of these maps. Closure of $\Sigma$ under composition suffices when type maps lengthen a run. The weak-vertical condition is preserved because every component of the concatenation was a component of an input.

The nerve preserves products as an equality of simplicial sets. The functor $\operatorname{Ex}$ preserves products as a right adjoint, and filtered colimits of sets commute with finite products. Therefore $\operatorname{Ex}^{\infty}$ preserves finite products and the terminal simplicial set. Apply it to the displayed multiplication and its commuting associativity and unit diagrams. Its usual natural weak equivalence from the identity and its Kan property give the asserted enrichment. For completeness, its Segal nerve in degree $n$ is
\[
 \coprod_{a_0,\ldots,a_n}\prod_{j=1}^{n}
 \operatorname{Ex}^{\infty}\mathcal G(a_{j-1},a_j).
\]
The Segal maps are isomorphisms. A Reedy fibrant replacement gives a Segal space, since the object space is discrete and the displayed products already compute the relevant homotopy fibre products. Rezk completion changes the presentation, not its associated enriched homotopy theory. These observations establish coherence for the raw model; a comparison must additionally identify its mapping spaces and its action on objects.
\end{proof}

\begin{result}{C.2}{Known}{The hammock enrichment has the localization interpretation}
For every small relative category $(M,\Sigma)$, the simplicial category $L^H(M,\Sigma)$ and its mappingwise $\operatorname{Ex}^{\infty}$ replacement have strict associative unital composition, component category $M[\Sigma^{-1}]$, and the localization mapping spaces. Composition with a specified arrow of $\Sigma$ induces weak equivalences on all relevant mapping spaces. These assertions hold without two-out-of-three or saturation.
\end{result}
\begin{proof}[Verification of the source and its scope]
The construction and composition are Definition 2.1 of \citet{DwyerKan1980Calculating}; Proposition 3.1 identifies components and Proposition 3.3 treats composition with a weak equivalence. Part B specifies the Grothendieck-construction version of the weak-vertical model and its comparison. The product-preserving Kan replacement is justified in C.1. None of these theorems asserts that an inclusion permitting arbitrary vertical components is a weak equivalence. Nor does an equivalence of individual mapping spaces, without compatibility, identify a prescribed multiplication. The all-arrows compatibility supplied in CATHOT E.24 and the companion's VII ms.~169 block is used with that precise restriction.
\end{proof}

\subsection{Pointed composition and Quillen's groupoid}

\begin{result}{C.3}{Proved}{The unit for pointed composition is $S^0$}
Composition of pointed mapping spaces, with zero morphisms as basepoints, has the form
\[
 K(b,c)\wedge K(a,b)\longrightarrow K(a,c),\qquad
 S^0\longrightarrow K(a,a).
\]
The nonbasepoint of the unit is $1_a$. Replacing smash product by the Cartesian monoidal structure on pointed spaces cannot express unital composition with absorbing zero whenever $1_a\ne0$ in $\pi_0K(a,a)$.
\end{result}
\begin{proof}
In the Cartesian monoidal category of pointed spaces the unit is the one-point space, and its only pointed map to $K(a,a)$ selects zero. The left unit axiom would then say $0\circ f=f$, while absorption says $0\circ f=0$. Applying this to $f=1_a$ contradicts the hypothesis on components. For smash product the unit is $S^0$, and a based map from it can select the identity at its other point. A map on the product which is constantly zero on either axis descends to the smash product. The associative and unit diagrams descend as well. This is the correction made in IX ms.~17; it is a structural requirement, not a choice of notation.
\end{proof}

\begin{result}{C.4}{Proved}{A sufficient pointed model with strict absorption}
Let $M$ be a pointed simplicial model category in a universe in which the objects and replacements under consideration exist. On its full simplicial subcategory of cofibrant-fibrant objects, the simplicial mapping spaces are pointed Kan complexes, their composition factors through smash product, and their units and associativity are strict. They present the pointed derived mapping spaces of $M$.
\end{result}
\begin{proof}
Let $0$ be the zero object. Tensoring preserves the initial object, and mapping into the terminal object is unique, so the enriched spaces $\Map_M(0,b)$ and $\Map_M(a,0)$ are each the one-point simplicial set. Their composite selects a zero simplex in $\Map_M(a,b)$ and all its degeneracies. Enriched associativity implies that composing with it in either variable gives that zero map exactly. Thus composition descends levelwise through the quotient by the two axes, which is the smash product of pointed simplicial sets. The identity maps define units from $S^0$. For cofibrant $a$ and fibrant $b$, the simplicial model axiom gives the Kan property. The derived identification is \citet[Corollary 4.7]{DwyerKan1980Function}; it is compatible with composition on the cofibrant-fibrant simplicial category, as in their Proposition 4.8. Choosing one cofibrant-fibrant representative for each desired object gives the stated model. This construction fixes a pointed enrichment, rather than choosing unrelated points in spaces already identified up to unpointed equivalence.
\end{proof}

\begin{result}{C.5}{Known}{The precise groupoid comparison with Quillen}
In a model category, with cofibrant source $a$ and fibrant target $b$, Quillen's groupoid of maps and homotopy classes of homotopies is naturally equivalent to the fundamental groupoid of a derived function complex. In a simplicial model category this is the fundamental groupoid of $\Map_M(a,b)$. In particular, for actual maps $f,g:a\to b$,
\[
 \pi_1^{\mathrm{Quillen}}(a,b;f,g)
 \cong\Hom_{\Pi_1\Map_M(a,b)}(f,g).
\]
In the pointed case the group at the zero map is $\pi_1(\Map_M(a,b),0)$. For arbitrary source and target the formula is applied after replacements, with specified transported maps.
\end{result}
\begin{proof}[Source check and transport]
\citet{Quillen1967} constructs the groupoid in Chapter I, \S2, Proposition 1. In the attached transcription this is Proposition 1.2.1. Chapter II, \S2, the remark following Proposition 5 explicitly identifies it with the fundamental groupoid of the simplicial mapping complex when the source is cofibrant and the target fibrant; the transcription numbers these as Proposition 2.2.5 and its following remark. For a general model category, the resolution description of function complexes in \citet[Proposition 4.4]{DwyerKan1980Function} realizes the same cylinder and path homotopies, and their relative homotopies. Its fundamental groupoid is consequently the same one.

A weak equivalence of mapping spaces induces an equivalence of fundamental groupoids, not equality of their chosen objects. To obtain the displayed bijection with the specified $f,g$, use their vertices under the comparison, or supply paths from their images to the selected vertices. Replacing a chosen basepoint path changes an isomorphism of fundamental groups by the corresponding conjugation. The fundamental groupoid is the formulation free of that ambiguity. Every component of the hammock complex already has zigzag vertices; its fundamental groupoid therefore extends the construction to derived morphisms. For a literal Grothendieck path vertex, the extension requires the comparison in Part B and its agreement with that vertex. This qualification is the remaining model-specific issue in IX-040, IX-041 and IX-072.
\end{proof}

\begin{result}{C.6}{Proved}{Loop and suspension identities with their actual hypothesis}
Let $\mathcal C$ be a pointed $\infty$-category. If the pullback $\Omega b=0\times_b0$ exists, then naturally in $a$,
\[
 \Map_{\mathcal C}(a,\Omega b)\simeq
       \Omega\Map_{\mathcal C}(a,b).
\]
If the pushout $\Sigma a=0\amalg_a0$ exists, then naturally in $b$,
\[
 \Map_{\mathcal C}(\Sigma a,b)\simeq
       \Omega\Map_{\mathcal C}(a,b).
\]
These are pointed equivalences with the zero maps as basepoints. When both constructions exist for every object, they give $\Sigma\dashv\Omega$. The corresponding construction in a pointed derivator gives an adjunction of derivator morphisms; $\Omega$ preserves homotopy limits and $\Sigma$ preserves homotopy colimits in its domain of shapes.
\end{result}
\begin{proof}
The defining mapping-space property of the first pullback is
\[
 \Map(a,\Omega b)\simeq
 \Map(a,0)\times^h_{\Map(a,b)}\Map(a,0).
\]
Both outside spaces are contractible, and their maps select zero. The right side is precisely the based loop space. The equivalence carries the zero map to its constant loop, since it is induced by the zero cone. The pushout calculation is the dual calculation in the first variable. Combining the two natural equivalences gives the adjunction, including its coherence. On components it yields
\[
 \Hom_{\Ho\mathcal C}(\Sigma a,b)
 \cong\Hom_{\Ho\mathcal C}(a,\Omega b)
 \cong\pi_1(\Map_{\mathcal C}(a,b),0),
\]
with the group operations induced by concatenation of loops.

For a pointed derivator, extension by zero to a corner, followed by the appropriate homotopy Kan extension to a square and evaluation at its opposite corner, constructs these functors simultaneously on all diagram categories. Their adjunction is \citet[Definition 3.16 and Proposition 3.17]{Groth2013}. The adjoint-morphism criterion and its Kan-extension compatibility in that paper, Proposition 2.9 and its dual, show that the right adjoint preserves homotopy limits and the left adjoint preserves homotopy colimits. Thus IX-084 and IX-087 hold in this setting. Merely writing $A\mapsto M^A[W_A^{-1}]$ has not established that setting.
\end{proof}

\begin{result}{C.7}{Refuted}{Representing the groups $\pi_1$ does not give the loop-space identity}
Even for a pointed $\infty$-category, representability of all functors $\pi_1\Map(a,-)$ and $\pi_1\Map(-,b)$ on its homotopy category does not imply the space-level identities of C.6.
\end{result}
\begin{proof}
Take the full subcategory of pointed spaces on $*,S^0,K(\mathbb Z,2)$. It is pointed. All of its mapping spaces have trivial fundamental group at zero. Indeed the only relevant nondiscrete one between distinct nonzero objects is
\[
 \Map_*(S^0,K(\mathbb Z,2))\simeq K(\mathbb Z,2).
\]
Maps from the connected space $K(\mathbb Z,2)$ to the discrete pointed $S^0$ form a point. The higher homotopy groups of the self-mapping space of $K(\mathbb Z,2)$ are its reduced cohomology groups in degrees $2-i$; these vanish for $i\geq1$. The self-mapping space of $S^0$ is discrete. Maps involving $*$ are contractible. Consequently the two $\pi_1$ functors are the constant trivial-group functors and are represented by $*$; one may take both proposed functors on the homotopy category to be the constant-zero functor.

For $a=S^0$ and $b=K(\mathbb Z,2)$ the proposed identity would identify the contractible space $\Map_*(a,*)$ with $K(\mathbb Z,1)$. This is false. The example refutes the inference from the representability of groups, which is exactly what precedes IX ms.~23. It does not claim that a true homotopy pullback has the wrong universal property. The missing datum is the pullback or an equivalent coherent representing property.
\end{proof}

\subsection{The linear statement with a specified enhancement}

\begin{result}{C.8}{Proved}{The requested linear objects for complexes}
Let $k$ be a commutative ring and let $M=\operatorname{Ch}(k)$, with quasi-isomorphisms and cohomological grading. For complexes $a,b$ set
\[
 K_k(a,b)=\tau_{\leq0}\RHom_k(a,b)\quad\in D^{\leq0}(k).
\]
There are associative unital composition maps
\[
 K_k(b,c)\otimes_k^{\mathbf L}K_k(a,b)\longrightarrow K_k(a,c),
 \qquad k\longrightarrow K_k(a,a).
\]
Under Dold--Kan, the underlying pointed space is the derived mapping space of complexes, and
\[
 \pi_i\Map_{L(M,W)}(a,b)=H^{-i}\RHom_k(a,b)
       =\Hom_{D(k)}(a,b[-i])\quad(i\geq0),
\]
where $i=0$ denotes the underlying set of the indicated abelian group. The same assertions hold in any specified dg model whose derived mapping complexes and composition are given by such an enrichment.
\end{result}
\begin{proof}
Choose a dg-projective replacement $Qa\to a$ for every object, in the working universe. Such replacements and the projective model structure are verified in \citet[Proposition 3.1 and Theorem 3.2]{Keller2006DG}. The complexes $\Hom_k(Qa,Qb)$ compute derived Hom and have strictly associative dg composition. For a cochain complex $C$, good truncation has $C^n$ for $n<0$, $\ker(d:C^0\to C^1)$ in degree zero, and zero in positive degrees. Composition restricts to these truncations: negative degrees add to a nonpositive degree and the product of degree-zero cycles is a cycle. The Leibniz identity verifies that the restriction is a chain map, and each $1_{Qa}$ is a degree-zero cycle.

For the derived tensor assertion we must not assume that these Hom complexes are $K$-flat. In the closed monoidal derived category, evaluation and its adjunction give
\[
 \RHom(b,c)\otimes_k^{\mathbf L}\RHom(a,b)\longrightarrow\RHom(a,c).
\]
The inclusion $D^{\leq0}(k)\to D(k)$ is monoidal: derived tensor of two nonpositive objects is nonpositive. Its right adjoint $\tau_{\leq0}$ is therefore lax monoidal. Apply that structure and the displayed evaluation composition. This constructs the required maps with their associative and unit coherence. Equivalently, the map is the unique factorization through the good truncation of the composition after the two truncation inclusions; its source is nonpositive. Repeating this for any number of factors identifies both bracketings with the same iterated evaluation map.

Reindex the good truncation as a nonnegative chain complex by $C_i=C^{-i}$, and apply the inverse Dold--Kan functor of \citet[Theorem 14.24.3]{StacksDoldKan}. Its homotopy groups are $H_i(C)=H^{-i}(C)$. The dg function complex, or a simplicial resolution of the complexes, identifies this simplicial module with the derived function complex. That is the enrichment version of the resolution comparison in \citet[Proposition 4.4]{DwyerKan1980Function}. The zero chain map supplies the basepoint; the dg composition supplies the bilinear pairing. Reversing the order of the displayed tensor factors requires the Koszul symmetry, not an unsigned exchange of homogeneous maps. This proves the claimed refinement of IX-061, IX-065 and IX-070 in the stated dg setting.

For example, $a=k[0]$ and $b=k[1]$ give $K_k(a,b)=k[1]$, concentrated in cohomological degree $-1$, with mapping space $K(k,1)$. Thus the higher information is present even though the group of degree-zero derived morphisms is zero. In this model $\Sigma a=a[1]$ and $\Omega b=b[-1]$ agree with C.6.
\end{proof}

\begin{result}{C.9}{Open}{The manuscript's unaxiomatized additive generality}
For an arbitrary additive $k$-linear category $M$ and a wide class $\Sigma$ closed under finite direct sums, we do not establish that the literal categories $\Ch_\Sigma(M;a,b)$ admit composition-compatible refinements in $D^{\leq0}(k)$ agreeing with their nerves. Adding two-out-of-three, saturation, or an abstract suspension--loop equivalence on the homotopy category is not proved here to suffice. The exact remaining question is existence of the linear enhancement together with a comparison to those specified path categories.
\end{result}
\begin{proof}[Scope and failed reduction]
At IX ms.~18 Grothendieck explicitly asks for additional compatibility of $\Sigma$ with addition, and at ms.~19 suggests an equivalence between suspension and loops. These sentences do not specify an axiom system. C.8 proves a precise instance, and it must not replace the unspecified class by all dg model categories without recording the restriction. An additive structure on a homotopy category alone does not construct chain-valued composition or a comparison of its underlying spaces. A stronger established criterion is \citet[Theorem 1.3 and Definition 6.2]{DuggerShipley2007}: a stable, additive, combinatorial model category has enrichment in spectra of simplicial abelian groups. Their definition includes compatible additive cosimplicial resolutions (Definition 6.1), not merely additive Hom sets. Their chain-enrichment results are in \S9. This supplies another sufficient enhanced context; it does not verify that an arbitrary relative additive category satisfies those hypotheses. We searched the literature on additive simplicial localization, dg localization, and additive model enrichments. No argument closing the general comparison gap is claimed.
\end{proof}

\begin{result}{C.10}{Open}{Literal pointed and ind-pointed transfer}
For a general pointed relative category, the claims that the manuscript's specified stabilized-path pairings define an associative unital smash enrichment, and that its stage system defines the proposed pointed objects of $\operatorname{Ind}(\Hot_*)$ and $\operatorname{Ind}(\Cat\,\mathrm{hot}_*)$, remain conditional on a pointed, composition-compatible comparison or a direct construction. An unpointed mapping-space equivalence alone does not settle IX-044, IX-049 or IX-052 at this literal level.
\end{result}
\begin{proof}[The exact outstanding coherence]
A chosen zero-arrow vertex gives a pointed nerve, but concatenating it with a nonempty path generally gives a different path object. To factor through smash one needs absorption homotopies in both variables which agree at the two-axis intersection; to assert associativity one needs their compatibility with threefold concatenation and units. For an ind-object one additionally needs these choices compatible with every transition map. C.4 supplies all of this in a strict pointed simplicial model, and C.6 supplies the resulting loop identities. Transferring it to the exact categorical system requires specified comparison maps with these coherences. A bijection on components, or an isomorphism of unpointed ind-objects, provides none of them. The elementary cases $\Sigma=\operatorname{Iso}(M)$ and a chosen dg model have the evident discrete or enhanced pointed refinement; the general literal statement remains on the register.
\end{proof}
